\section{Introduction}\label{g:introduction}

Anticanonical nonvanishing asks when positivity of $-K_X$ produces a
nonzero section of $-mK_X$ for some integer $m>0$. We study metric
descent and dimension reduction for this problem, including sections
with prescribed poles and sections fixed by a torus.

Throughout the paper varieties are over $\mathbb C$. We write line
bundles additively. On a smooth variety $X$, the anticanonical bundle is
$-K_X=\det T_X$, with its natural linearization whenever a group acts.
A singular Hermitian metric is semipositive if its local squared frame
norm is $e^{-\varphi}$ for a plurisubharmonic weight $\varphi$.
A metric with locally bounded weights satisfies two-sided local bounds
on $\varphi$. Such a condition is stronger than pseudoeffectivity and
is preserved by the normalized maximum construction below.

Two section theorems are applications of the constructions developed
here. Let $X$ be smooth projective with no positive-degree holomorphic
forms, and let $D$ be an effective integral divisor. View its canonical
section $s_D$ as a $(-K_X+D)$-valued top form. If a semipositive metric
$h$ on $-K_X+D$ makes the measure $|s_D|_h^2$ finite, then for every
line bundle $L$ with locally bounded semipositive weights,
Theorem~\ref{g:allowed} gives integers $a>0$ and $c\ge0$ such that
\[
 H^0(X,aL+cD)\ne0.
\]
If $-K_X$ is smoothly semipositive and an algebraic torus $T$ acts on
$X$, Theorem~\ref{rb:invariant} gives
$H^0(X,-mK_X)^T\ne0$ for some $m>0$, for the natural action.
These conclusions are also established in the companion
\cite[Theorems~1.2 and~8.1]{HFiniteVolume}, which uses invariant sections
and finite-cover descent to obtain anticanonical nonvanishing without
the no-forms hypothesis \cite[Theorem~1.1]{HFiniteVolume}.
Here the emphasis is on reusable transfer statements, rank-preserving
contractions, and alternative constructions by orbit integrals and
finite-flat norms.

A section constructed on the base of a rational fibration need not
return to the original variety. Multiplying its pullback by a rational
relative section can leave poles over degenerate fibers, and a second
contraction can destroy the one-dimensional space that made an
integral metric positive. We study these two obstructions together:
the boundary corrections must disappear under birational pushdown,
and the contraction must preserve both a generic section and its
adjoint section ranks.

\subsection{From a volume on the total space to a metric on the base}

Let $X$ be smooth projective and suppose $L=-K_X$ has a smooth
semipositive metric. The metric determines a smooth positive volume
form on $X$. Resolve a rational fibration to obtain
\[
 X\xleftarrow{\ \pi\ } Z\xrightarrow{\ g\ }Y,
 \qquad E=K_{Z/X},
\]
with $X,Z,Y$ smooth projective. The adjoint identity is
\[K_{Z/Y}+\pi^*(-K_X)=E+g^*(-K_Y).\]
Over a smooth fiber, the Jacobian section
$s_E$ is an adjoint section after a base canonical frame has been
chosen. If it spans the generic adjoint space, integration along the
fibers turns the pushforward volume into the squared norm of a line.
Relative direct-image positivity then supplies the curvature inequality
on the smooth locus. The question is what happens at its boundary.

The first metric result, Theorem~\ref{g:toroidal-volume}, answers this question
on a normal equidimensional toroidal intermediate model
$W\to Y$ dominating $X$. For a base prime $Q$, let $e_D$ be its pullback
multiplicity at a prime $D$ of $W$ above $Q$, and let $a_D$ be the
Jacobian order over $X$. The correction is the explicit rational divisor
\[
 A_0=\sum_Q
 \max\left\{0,\min_{D\mapsto Q}
          \left(\frac{1+a_D}{e_D}-1\right)\right\}Q.
\]
Its pullback to $W$ is exceptional over $X$. The resulting metric on
$-K_Y+A_0$ has locally integrable squared frame norms without inserting
a vanishing section. If, away from the Jacobian divisor, the coefficient
curvature dominates a positive pulled-back base form on one patch,
then this metric is smooth and strictly positive on a nonempty ordinary
open set. The proof checks every divisorial valuation, including those
centered in intersections of the boundary. This construction is proved
before any section-existence argument. The companion
\cite[Sections~5.3 and~11.2]{EControlledBoundary} uses precisely this
unweighted integrability, strictness, and exceptional support for its
fixed-error effectivity argument and its Fano-type model of an auxiliary
base.

The distinction between a normal intermediate model and its smooth
resolution is useful. Equidimensionality controls which divisors can
carry poles. Smoothness permits fiber integration. A smooth resolution
need not be equidimensional, but its additional divisors over subsets of
base codimension at least two are exceptional over the original variety.
Keeping both models makes the final support statement exact.

\subsection{A contraction that keeps its generic section}

The second main result is Theorem~\ref{g:cf:contraction-theorem}.
Let $X$ have no positive-degree holomorphic forms, let $-K_X$ be smoothly
semipositive, and let a torus $T$ act on $X$. Consider an invariant
rational fibration to a trivially acted-on base $Y$ whose function field is relatively
algebraically closed in $\mathbb C(X)$. Assume that a positive
anticanonical multiple has a nonzero invariant section on the generic
fiber, and put $\eta=\operatorname{Spec}\mathbb C(Y)$. Suppose that the invariant generic adjoint spaces
\[
 H^0\bigl(Z_\eta,(K_{Z/X}+m\pi^*(-K_X))|_{Z_\eta}\bigr)^T
\]
have dimension one for every positive multiple of a fixed integer.
Then either a positive anticanonical multiple has a global invariant
section, or the fibration factors rationally through a strictly
lower-dimensional base satisfying the same two conditions.

The two assumptions have different roles. The generic section supplies
the line whose high moments yield a bounded base metric. Multiplication
by its powers also compares degree zero with the positive-degree
adjoint spaces and forces rank one at zero; this supplies the zeroth
density as well as the high moments. The rank condition makes ratios of sections into base
functions. After contraction, the proof must recover that statement
for the larger generic fiber, as well as producing its section.

Here is the mechanism. Vertical minima normalize the section to a
regular map from a base line. A toroidal density calculation produces
both a bounded metric on that line and a nef anti-log-canonical class
with a signed boundary. The negative boundary is supported on
those base primes where section poles can be absorbed upstairs.
If a suitable small boundary correction is effective, its section
lifts to $X$. Otherwise a movable class annihilates the relevant
divisors. Twisted differential forms give an unbounded sequence of
effective divisors of degree zero. Their ratios define a smaller
field. Affine relations among their horizontal multiplicities give the
new generic section, and the controlled base poles restore all the
adjoint ranks. This last step is what permits iteration.

\subsection{Methods, predecessors, and other metric hypotheses}

The distinction between numerical and linear effectivity is important
in positioning these results. Lazi\'c--Matsumura--Peternell--Tsakanikas--Xie
prove numerical effectivity of nef anti-log-canonical divisors for projective
log canonical threefold pairs under $\mathbb Q$-factoriality or rational
singularities \cite[Theorem~A]{LMPTX2023}. A numerical equivalence with
an effective divisor need not supply a section of a multiple of the
given line. M\"uller proves actual nonvanishing for projective klt
threefold pairs with nef anti-log-canonical divisor, and in arbitrary
dimension when that divisor is semiample on a general
maximal-rationally-connected fiber
\cite[Theorem~A and Corollary~B]{Muller2025}.
His equivariant theorem gives invariant plurisections for commutative
linear algebraic groups on projective sub-log-canonical pairs with
semiample anti-log-canonical divisor \cite[Theorem~C]{Muller2025}.
Our invariant conclusion instead uses smooth semipositivity on a
smooth variety without positive-degree forms and treats algebraic
tori. It does not assert the same generality of pairs or groups.
M\"uller's moment-polytope argument also precedes the toric initial-section
construction in Section~\ref{g:cf:fields}
\cite[Proposition~3.13 and Corollary~3.14]{Muller2025}.

The birational preparations use flattening by Raynaud--Gruson
\cite{RaynaudGruson1971}, the toroidalization of
Abramovich--Denef--Karu \cite{ADK2013}, and the cone subdivisions and
weak semistable reduction of Abramovich--Karu \cite{AK2000}.
Birational toroidal preparation and reduced fibers after alteration are
different constructions here. We use the former for valuation and
support estimates; the latter gives a separate norm construction in
Section~\ref{g:reduced-flat}.

The curvature mechanism builds on Berndtsson's smooth direct-image
positivity \cite[Theorem~1.2]{Berndtsson2009} and the relative Bergman
theorem of Berndtsson--P\u aun \cite[Theorem~0.1]{BerndtssonPaun2008}.
The singular direct-image formulation is due to P\u aun--Takayama
\cite[Theorem~3.3.5]{PT2018}; the actual-multiplier-ideal formulation is
also presented in \cite[Theorem~21.1]{HPS}. In the ordinary singular case
we verify integrability of the generating adjoint section, including
the actual multiplier ideal. In the invariant case smooth coefficients
allow compact averaging to give an orthogonal holomorphic quotient.
This explains why the singular and invariant forms of the transfer
statement have separate hypotheses. The normalized high-moment limit
then produces a bounded metric, while the zeroth moment retains a
finite-volume donor metric on a different line bundle.

Movable-class slopes connect those metrics to dimension reduction.
We use divisor--curve duality \cite{BDPP2013}, movable slope theory
\cite{GKP2016}, and the algebraicity criterion of Campana--P\u aun
\cite{CampanaPaun2019}. The foliation argument follows the strategy
used by Ou for generic nefness \cite{Ou2023}; the boundary and
integrability conditions required here are retained in each local
version. Proposition~\ref{g:bt:weighted-slope} proves the general finite-volume
cotangent estimate by approximating a rational boundary while preserving
its strict discrepancy inequalities. For integral boundaries, the companion
\cite[Theorem ``Finite-volume cotangent bound'']{HFiniteVolume} gives a
different proof through ramification estimates and relative Bergman
metrics. Hard Lefschetz with multiplier ideals
\cite[Theorem~0.1]{DPS2001} supplies twisted forms. Their determinant
extraction follows Lazi\'c--Peternell \cite[Lemma~4.1]{LP2018}, as adapted
in \cite[Lemma~5.1]{LMPTX2023}. The horizontal divisor comparison is
a relative version of the three-index argument in
\cite[Theorem~3.1, Step~3]{LP2018}: ratios in the smaller base field
give equality of horizontal parts instead of the absolute divisor
equality furnished there by Iitaka dimension zero. The subsequent
control of vertical poles and restoration of adjoint rank are separate
steps.

Several useful variants require different input data. An integrable
metric after multiplication by a boundary section need not have
integrable frame norms. A signed sub-klt divisor is also different
from an effective klt boundary. We therefore distinguish weighted,
unweighted, and signed-boundary transfer statements. The large-moment
construction uses only eventual adjoint ranks in its curvature step.
Together with its regular distinguished section those ranks imply
rank one also in the earlier degrees, but the construction does not
use zeroth-moment curvature. Each variant is carried
through its support and rank verification before it meets a common
contraction or induction.

Two further constructions illuminate the metric step. On a dense torus
orbit, logarithmic coordinates turn the volume density into an integral
of a log-concave function. This is the setting of complex Pr\'ekopa
positivity and Kiselman's minimum principle
\cite{Kiselman1978,Berndtsson1998,Berndtsson2006}.
The Brascamp--Lieb variance calculation
\cite{BrascampLieb1976,KolesnikovMilman2017} proves that the negative logarithms of its integral and its supremum
are plurisubharmonic in the base variables; Section~\ref{g:orbits} gives the complete
argument and its invariant-section application. On a reduced flat
alteration, a generic section instead determines a line nonzero on
every fiber. Its finite-flat determinant norm returns that line to a
birational model of the original base. The norm multiplies the section
exponent by its degree and preserves the required polar support.

\subsection{Reading the proofs}

Section~\ref{g:sec-toroidal-volume} proves the exceptional unweighted
metric construction. Section~\ref{g:analytic} gives two-metric transfer
and its rank-one metric hypotheses; Section~\ref{g:cf:contraction} then
proves the rank-preserving contraction. Section~\ref{g:bt} establishes
the weighted, unweighted, and signed-boundary variants.

For the ordinary allowed-pole theorem, the direct route is from the
two-metric transfer in Section~\ref{g:analytic}, through the weighted
cotangent estimate, to the induction in Section~\ref{g:reductions}.
For invariant sections, Section~\ref{g:orbits} applies that theorem
using either rank-one quotient integrals or explicit logarithmic orbit
integration. These proofs do not require the rank-preserving
contraction theorem. Sections~\ref{g:cf:fields}
and~\ref{g:field-variants} then identify the ratio fields, extract their
generic sections, and vary the base construction. The boundary
applications in Section~\ref{g:ba} and multiplicative systems in
Section~\ref{g:multiplicative} give other contractions with their own
rank hypotheses.
Section~\ref{g:reduced-flat} replaces the metric-transfer step by an
alteration and a finite-flat norm, then proves the rationally connected
cases of the ordinary and invariant conclusions over the original base
field.
